%-----------------------------------------------------------------------------%
\chapter{\babSatu}
\label{bab:1}
%-----------------------------------------------------------------------------%
Pada bab ini, akan dijelaskan tentang latar belakang dan permasalahan yang diselesaikan pada penelitian ini.

%-----------------------------------------------------------------------------%
\section{Deskripsi Masalah}
\label{sec:deskripsiMasalah}
%-----------------------------------------------------------------------------%
\todo{Sebutkan permasalahan penelitian Anda dari latar belakang tersebut.}

%-----------------------------------------------------------------------------%
Fungsi parametrik $(x(t),y(t))$ dapat merepresentasikan gerak suatu objek pada bidang dua dimensi. Fungsi tersebut menghasilkan koordinat objek pada detik ke-$t$.
Jika diberikan $(n+1)$ titik $(x_0,y_0),(x_1,y_1),(x_2,y_2),...,(x_n,y_n)$ dengan asumsi sebagai berikut:
\begin{itemize}
	\item $x_i \neq x_j$ dan $y_i \neq y_j$
	\item $(x_i,y_i)$ adalah posisi objek yang terdeteksi dan terekam pada detik kei$i$.
\end{itemize}
Dengan menggunakan \f{spline interpolation} varian \f{natural cubic spline}, \f{deliverables} yang diharapkan antara lain:
\begin{itemize}
	\item Kurva \f{smooth} yang menggambarkan perkiraan jalur yang dilalui oleh objek
	\item Grafik jalur jalur yang ditempuh oleh objek
\end{itemize}

% \subsection{Definisi Permasalahan}
% \label{sec:definisiMasalah}
% %-----------------------------------------------------------------------------%
% Berikut ini adalah rumusan permasalahan dari penelitian yang dilakukan:
% \begin{itemize}
% 	\item Bagaimana cara membuat pertanyaan penelitian?
% \end{itemize}
% \todo{Tuliskan permasalahan yang ingin diselesaikan. Bisa juga berbentuk pertanyaan}


% %-----------------------------------------------------------------------------%
% \subsection{Batasan Permasalahan}
% \label{sec:batasanMasalah}
% %-----------------------------------------------------------------------------%
% Berikut ini adalah asumsi yang digunakan sebagai batasan penelitian ini:
% \begin{itemize}
% 	\item Salah satu batasannya adalah, ini hanya \f{template}.
% \end{itemize}
% \todo{Umumnya ada asumsi atau batasan yang digunakan untuk menjawab pertanyaan-pertanyaan penelitian diatas.}


% %-----------------------------------------------------------------------------%
% \section{Tujuan Penelitian}
% \label{sec:tujuan}
% %-----------------------------------------------------------------------------%
% Berikut ini adalah tujuan penelitian yang dilakukan:
% \begin{itemize}
% 	\item Untuk memberikan \f{template} yang dapat mempermudah skripsi orang lain.
% \end{itemize}
% \todo{Tuliskan tujuan penelitian Anda di bagian ini.}


% %-----------------------------------------------------------------------------%
% \section{Posisi Penelitian}
% \label{sec:posisiPenelitian}
% %-----------------------------------------------------------------------------%
% \todo{Sebutkan posisi penelitian Anda. Ada baiknya jika Anda menggunakan gambar atau diagram. Template ini telah menyediakan contoh cara memasukkan gambar.}

% \begin{figure}
% 	\centering
% 	\includegraphics[width=0.4\textwidth]{assets/pics/makara.png}
% 	\caption{Penjelasan singkat terkait gambar.}
% 	\label{fig:research_position}
% \end{figure}

% \todo{Jelaskan \pic~\ref{fig:research_position} di sini.}


% %-----------------------------------------------------------------------------%
% \section{Langkah Penelitian}
% \label{sec:langkahPenelitian}
% %-----------------------------------------------------------------------------%
% Berikut ini adalah langkah penelitian yang telah dilakukan:
% \begin{enumerate}
% 	\item Tinjauan literatur \\
% 	      Pada tahap ini, dipelajari teori-teori yang terkait dengan penelitian ini untuk mendapatkan konsep dasar yang dibutuhkan dalam mencapai tujuan penelitian.
% 	\item Analisis implementasi dan kesimpulan \\
% 	      Pada tahap ini, digunakan studi kasus untuk analisis terkait kegunaan \f{template}. Setelah melakukan analisis tersebut, ditarik kesimpulan keseluruhan dari penelitian ini.
% \end{enumerate}


% %-----------------------------------------------------------------------------%
% \section{Sistematika Penulisan}
% \label{sec:sistematikaPenulisan}
% %-----------------------------------------------------------------------------%
% Sistematika penulisan laporan adalah sebagai berikut:
% \begin{itemize}
% 	\item Bab 1 \babSatu \\
% 	      Bab ini mencakup latar belakang, cakupan penelitian, dan pendefinisian masalah.
% 	\item Bab 2 \babDua \\
% 	      Bab ini mencakup pemaparan terminologi dan teori yang terkait dengan penelitian berdasarkan hasil tinjauan pustaka yang telah digunakan, sekaligus memperlihatkan kaitan teori dengan penelitian.
% 	\item Bab 3 \babTiga \\
% 	      Apa itu Bab 3?
% 	\item Bab 4 \babEmpat \\
% 	      Apa itu Bab 4?
% 	\item Bab 5 \babLima \\
% 	      Apa itu Bab 5?
% 	\item Bab 6 \kesimpulan \\
% 	      Bab ini mencakup kesimpulan akhir penelitian dan saran untuk pengembangan berikutnya.
% \end{itemize}

% \todo{Anda bisa mengubah atau menambahkan penjelasan singkat mengenai isi masing-masing bab. Setiap tugas akhir pasti ada yang berbeda pada bagian ini.}
